% ===========================================================================
%  Capítulo 5 — Entre Aventuras
% ===========================================================================
\chapter{Entre Aventuras}

Nem toda sessão é masmorra. Este capítulo reúne os quatro sistemas que rodam entre combates
--- tempo livre, fama na Guilda, reputação com o mundo e o que dá pra fabricar com as
próprias mãos --- porque eles se usam com a mesma frequência que qualquer regra do Capítulo
4, mesmo fora da luta.

% ---------------------------------------------------------------------------
\section{Tempo Livre e Downtime}

Downtime é contado em \textbf{blocos de 1 semana} --- a mesma unidade que o repouso longo do
Capítulo 4 já usa pra curar todos os PV. No fim de cada semana livre, cada personagem
escolhe \textbf{uma} atividade da lista abaixo.

\begin{tabelalivro}{L{4cm} L{9.8cm}}{Atividades de Downtime}
\cabtabela{\textbf{Atividade} & \textbf{Efeito}}
Treinar & Ganhe Vantagem no próximo teste de uma Perícia à escolha, ligada à sua Árvore Inicial ou a uma Perícia que você já tenha --- dura até ser usado ou até 1 mês passar. Não concede PA. \\
Recuperar-se & Como o repouso longo de uma semana: todos os PV são restaurados, e mais 1 nível de Exaustão é removido além do normal. \\
Trabalhar & Ganhe PO igual a 2d6 $\times$ seu maior Bônus de Rank (mínimo 2d6), pelo seu Ofício, sua fama ou um trabalho comum da cidade. \\
Cultivar um Contato & Anote um NPC nomeado e uma cidade ou facção. Da próxima vez que precisar de uma informação ou favor pequeno, o Mestre pode deixar esse contato resolver --- sem PP, sem teste. \\
Estudar um Ofício ou Ritual & Com a Perícia de Ofícios ligada ao que quer fazer, produza um item mundano ou prepare os materiais de um ritual que já pode conjurar. \\
Vigiar as Costas do Grupo & Sem efeito próprio, mas concede a outro personagem Vantagem na atividade dele nesta semana. \\
\end{tabelalivro}

\begin{aviso}{Downtime não compra progressão}
    Nenhuma atividade acima concede PA, magia, talento ou Rank --- isso só vem de jogar a
    campanha. Downtime existe pra que o tempo entre aventuras pareça vivido, não pra virar
    uma segunda forma de subir de patamar sem risco.
\end{aviso}

\begin{nota}{Downtime Interrompido}
    Se uma aflição (Cap. 4, §7) estiver ativa em alguém do grupo, a Profundidade dela
    continua subindo normalmente durante o downtime --- um bloco de \enquote{Recuperar-se}
    não pausa o relógio de um veneno ou de uma doença.
\end{nota}

% ---------------------------------------------------------------------------
\section{A Guilda de Aventureiros}

\subsection{O Rank Não É o Patamar}

\textbf{O Rank de Aventureiro (F a S) mede reputação, não poder de combate.} Ele não aparece
em nenhuma fórmula deste livro, não dá bônus de ataque, e não é igual ao Rank das suas
Árvores de Progressão. Um Deus da Espada desconhecido que nunca aceitou um contrato formal
pode ser Rank F. Um grupo de Rank A pode ter só um patamar Avançado cada --- a diferença é
que eles já resolveram cem contratos e a Guilda sabe o nome deles.

\begin{tabelalivro}{C{1.2cm} L{5.6cm} L{6.8cm}}{Ranks de Aventureiro}
\cabtabela{\textbf{Rank} & \textbf{Feito representativo} & \textbf{O que muda}}
F & Recém-registrado --- ainda não fez nada que a sede saiba. & Só pega contrato de mural público, sem escolta nem garantia. \\
E & Sobreviveu ao trabalho de rotina algumas vezes. & Escolta de caravana, extermínio de pragas, entrega em estrada segura. \\
D & Resolveu algo fora da cidade-sede sem apoio da Guilda. & Aceita contratos fora da cidade-sede. Pagamento sobe; a Guilda cobra 10\% de taxa. \\
C & Liderou outros aventureiros num contrato e todos voltaram. & Pode liderar um grupo de Ranks inferiores --- e responde por eles. Acesso ao arquivo de bestas da sede local. \\
B & Resolveu algo que chegou aos ouvidos de um nobre ou general. & Contratos de nobreza e de guerra pequena passam pela sua mesa. \\
A & Fez algo que virou boato em mais de uma cidade. & Reconhecido em qualquer continente com Guilda. Recusar um contrato regional exige justificativa formal. \\
S & Fez algo que devia ter sido impossível. & Menos de dez vivos por continente. Contratado direto por reinos e Guildas de outras nações. \\
\end{tabelalivro}

\begin{nota}{A tabela é exemplo, não fórmula}
    Não existe PA, teste ou compra que suba o Rank de Aventureiro. \textbf{O Mestre fala
    quando o personagem sobe}, olhando pro que o grupo \emph{fez} publicamente, não pro que
    gastou na ficha: um Deus da Espada que resolveu tudo em segredo pode morrer Rank F ---
    ele é forte, só não é famoso.
\end{nota}

\begin{aviso}{O palpite que a ficha mostra não é esta regra}
    A ficha digital exibe um Rank de Aventureiro estimado a partir do PA já gasto, porque uma
    ficha em branco precisa mostrar \emph{alguma coisa} naquele campo. \textbf{Esse número não
    é regra} --- ele contradiz tudo o que esta seção diz, e existe só como palpite inicial pra
    mesa que ainda não decidiu. Assim que o Mestre fixar o Rank, ele passa a valer e a
    estimativa some. Se a sua mesa quiser, ignore o campo inteiro.
\end{aviso}

\subsection{Subindo de Rank e Obrigações}

A promoção nunca é automática, mesmo depois de o Mestre decidir que o feito foi grande o
bastante: exige voltar à sede, ser avaliado, e --- a partir de Rank C --- pagar uma taxa de
registro em PO.

A partir de \textbf{Rank C}, recusar um contrato marcado como emergência sem justificativa
perde Rank. A partir de \textbf{Rank A}, a morte do aventureiro em contrato é investigada
formalmente pela sede.

\subsection{A Loja da Guilda}

Rank de Aventureiro não é só fama --- é a credencial que abre a porta do que a Guilda deixa
você comprar. Cada Rank libera o próximo andar do catálogo; abaixo do seu Rank, o item
simplesmente não está à venda ali, por mais PO que você tenha.

% GERADO — fonte: src/data/shopItems.ts
\begin{tabelalivro}{L{4.8cm} C{1.6cm} C{1.4cm} L{5.6cm}}{Armas --- Rank de Guilda é o mínimo pra o item aparecer à venda}
\cabtabela{\textbf{Item} & \textbf{Preço} & \textbf{Rank} & \textbf{Notas}}
Adaga / Punhal & 6 PO & F & Arma mundana, sem encantamento --- ver Cap. 3, "O Dado de Arma", pra como ela escala com seu Rank. \\
Funda / Dardo & 6 PO & F & Arma mundana, sem encantamento --- ver Cap. 3, "O Dado de Arma", pra como ela escala com seu Rank. \\
Chicote & 6 PO & F & Arma mundana, sem encantamento --- ver Cap. 3, "O Dado de Arma", pra como ela escala com seu Rank. \\
Espada Curta & 15 PO & F & Arma mundana, sem encantamento --- ver Cap. 3, "O Dado de Arma", pra como ela escala com seu Rank. \\
Objeto Improvisado & 15 PO & F & Arma mundana, sem encantamento --- ver Cap. 3, "O Dado de Arma", pra como ela escala com seu Rank. \\
Arco Curto & 15 PO & F & Arma mundana, sem encantamento --- ver Cap. 3, "O Dado de Arma", pra como ela escala com seu Rank. \\
Rapieira & 15 PO & F & Arma mundana, sem encantamento --- ver Cap. 3, "O Dado de Arma", pra como ela escala com seu Rank. \\
Espada Longa & 35 PO & F & Arma mundana, sem encantamento --- ver Cap. 3, "O Dado de Arma", pra como ela escala com seu Rank. \\
Machado de Batalha & 35 PO & F & Arma mundana, sem encantamento --- ver Cap. 3, "O Dado de Arma", pra como ela escala com seu Rank. \\
Arco Longo & 35 PO & F & Arma mundana, sem encantamento --- ver Cap. 3, "O Dado de Arma", pra como ela escala com seu Rank. \\
Foice de Guerra & 35 PO & F & Arma mundana, sem encantamento --- ver Cap. 3, "O Dado de Arma", pra como ela escala com seu Rank. \\
Espadão / Montante & 65 PO & F & Arma mundana, sem encantamento --- ver Cap. 3, "O Dado de Arma", pra como ela escala com seu Rank. \\
Martelo de Guerra & 65 PO & F & Arma mundana, sem encantamento --- ver Cap. 3, "O Dado de Arma", pra como ela escala com seu Rank. \\
Alabarda / Lança & 65 PO & F & Arma mundana, sem encantamento --- ver Cap. 3, "O Dado de Arma", pra como ela escala com seu Rank. \\
Besta & 65 PO & F & Arma mundana, sem encantamento --- ver Cap. 3, "O Dado de Arma", pra como ela escala com seu Rank. \\
Adaga de Prata & 40 PO & E & Liga de prata batida em vez de aço comum --- eficaz contra certas criaturas amaldiçoadas ou licantrópicas, a critério do Mestre. Sem bônus numérico contra o resto. \\
Lâmina Balanceada & 150 PO & C & Peso distribuído com perfeição rara --- funciona em qualquer arma leve, não só espadas. +1 no teste de acerto (nunca no dano). \\
Machado Sanguessedento & 350 PO & B & Ao reduzir um inimigo a 0 PV com este golpe (nunca em dano que não mata), você recupera 1d4 PV. \\
Lança Persecutora & 550 PO & A & Vantagem no teste de acerto contra um alvo que se afastou de você neste turno --- nunca perde o rastro de quem tenta fugir. \\
Espada-Fantasma & 800 PO & S & Uma vez por Descanso Longo, um ataque com ela ignora completamente a CA do alvo (acerto automático) --- o dano continua normal, sem bônus. \\
Espada Corta-Aço & 260 PO & D & Propriedade invertida: quanto mais duro o alvo, mais fácil ela acerta. Contra CA 15+, +1 no teste de acerto; CA 20+, +2; CA 25+, +3 (teto). O Dado de Arma não muda --- é só acerto, não dano extra (dano extra é Encantamento, Cap. 5 §4). Contra algo mole ou frágil sem resistência estrutural de verdade --- pano, corda, uma vela, uma folha de papel --- ela simplesmente não corta, não importa a força do golpe: desliza sem efeito. \\
\end{tabelalivro}

\begin{tabelalivro}{L{4.8cm} C{1.6cm} C{1.4cm} L{5.6cm}}{Armaduras --- Rank de Guilda é o mínimo pra o item aparecer à venda}
\cabtabela{\textbf{Item} & \textbf{Preço} & \textbf{Rank} & \textbf{Notas}}
Escudo & 25 PO & F & Empunhado numa mão só --- soma com qualquer outra armadura vestida. \\
Armadura Leve (couro) & 20 PO & F & Peito e braçadeiras de couro batido --- não atrapalha a mobilidade. \\
Armadura Média & 60 PO & F & Cota de malha com reforços --- o padrão de quem pretende segurar a linha de frente. \\
Armadura Pesada & 150 PO & E & Placas completas --- pesada o bastante pra sede exigir alguma reputação antes de vender. \\
Cota Élfica & 180 PO & D & Tecelagem que nenhum ferreiro comum reproduz --- tão leve quanto couro batido, protege quase como malha reforçada. \\
Armadura Rúnica & 260 PO & C & Runas gravadas na superfície distribuem o impacto de um golpe por toda a peça --- ligeiramente mais pesada que a Média, protege mais. \\
Placas Dracônicas & 450 PO & B & Escamas de dragão menor reforçando o metal --- Vantagem em testes de resistência contra Medo causado por criaturas dracônicas. Proteção física igual à Armadura Pesada. \\
Armadura Rúnica Maior & 650 PO & A & A evolução da Armadura Rúnica --- mais runas, mais peso, mais proteção. \\
Égide Lendária & 1200 PO & S & Uma vez por Descanso Longo, ignore completamente o dano de um único golpe recebido --- decida depois de ver o resultado do ataque. Se isso evitar que o dano leve você a 0 PV, conta como a sua Salvação do combate (Cap. 4, §4): não empilha com Aguentar, Rejeitar a Morte, Sem Baixas ou Custe o Que Custar. \\
\end{tabelalivro}

\begin{tabelalivro}{L{4.8cm} C{1.6cm} C{1.4cm} L{5.6cm}}{Equipamento de Aventura --- Rank de Guilda é o mínimo pra o item aparecer à venda}
\cabtabela{\textbf{Item} & \textbf{Preço} & \textbf{Rank} & \textbf{Notas}}
Corda (15m) & 2 PO & F & Trançada, aguenta o peso de uma pessoa em queda. \\
Tocha (feixe de 5) & 1 PO & F & Cerca de 1 hora de luz cada. \\
Kit de Primeiros Socorros & 10 PO & F & Bandagens, agulha, linha e um torniquete --- não substitui Cura, só estabiliza. \\
Ração de Viagem (1 semana) & 5 PO & F & Seca, pesada, dura em qualquer clima. \\
Cantil & 2 PO & F & Couro impermeabilizado, cerca de 1 litro. \\
Mochila de Couro & 8 PO & F & Reforçada, com alças duplas pra carregar mais sem cansar. \\
Kit de Acampamento & 15 PO & F & Tenda pequena, cobertor e um saco de dormir. \\
Isqueiro de Faísca & 3 PO & F & Pederneira e aço --- acende fogo mesmo com vento leve. \\
Lampião de Óleo & 12 PO & F & Luz mais estável que tocha, dura a noite toda com o tanque cheio. \\
Kit de Arrombamento & 18 PO & F & Gazuas e alavanca fina --- abre fechaduras e engates simples sem estragar a porta. \\
Gancho e Corda de Escalada & 10 PO & F & Corda de 20m com gancho de três pontas --- prende em beiradas, galhos ou grades. \\
Mapa em Branco & 6 PO & F & Pergaminho de boa gramatura, pena e tinta --- pra registrar o terreno que ninguém ainda desenhou. \\
Provisões de Longa Duração (30 dias) & 18 PO & F & A versão grande da Ração de Viagem --- mais pesada, mas rende um mês inteiro sem estragar. \\
Luneta & 30 PO & E & Lentes de vidro comum, sem magia --- enxerga detalhes a centenas de metros em dia claro. \\
Kit de Alpinismo Completo & 60 PO & D & Pitons, mosquetões e arnês --- escalada técnica em rocha viva, não só um gancho e uma corda. \\
Barraca Reforçada contra Clima Extremo & 90 PO & C & Aguenta nevasca, areia e vento forte sem rasgar --- o Kit de Acampamento comum não sobrevive nesses climas. \\
Provisões Concentradas de Campanha (90 dias) & 150 PO & B & Rações compactadas ao extremo --- três meses de comida no peso que a Ração de Viagem normal levaria pra duas semanas. \\
Kit de Sobrevivência de Expedição & 300 PO & A & Tudo que uma equipe precisa pra sobreviver semanas em terreno hostil desconhecido --- do tipo que só quem já liderou expedições sabe montar direito. \\
Mapa Master de um Continente Inteiro & 1000 PO & S & Levantamento cartográfico completo e atualizado --- a maioria dos mapas assim é segredo de estado. Um dos poucos itens mundanos que vale uma fortuna sozinho. \\
\end{tabelalivro}

\begin{tabelalivro}{L{4.8cm} C{1.6cm} C{1.4cm} L{5.6cm}}{Poções --- Rank de Guilda é o mínimo pra o item aparecer à venda}
\cabtabela{\textbf{Item} & \textbf{Preço} & \textbf{Rank} & \textbf{Notas}}
Poção Menor de Cura & 15 PO & F & Reproduz uma magia de Cura de rank Principiante ou Intermediário, sem precisar de mago presente. \\
Poção de Antídoto & 25 PO & E & Remove 1 ponto de Profundidade de uma única aflição (Cap. 4, §7). \\
Elixir de Foco & 40 PO & E & Vantagem no próximo teste de resistência de Espírito --- ajuda a resistir Trauma num momento específico. \\
Poção Maior de Cura & 60 PO & C & Reproduz uma magia de Cura de rank Avançado ou Santo. \\
Poção de Vigor Passageiro & 45 PO & D & Vantagem no próximo teste de resistência de Vigor --- a versão física do Elixir de Foco. \\
Poção Régia de Cura & 120 PO & B & Reproduz uma magia de Cura de rank Rei --- um degrau acima da Poção Maior. \\
Elixir de Regeneração & 200 PO & A & Remove toda a Exaustão acumulada de quem bebe (Cap. 4, §8) --- não cura PV nem PM, só o cansaço acumulado. \\
Poção Imperial de Cura & 400 PO & S & Reproduz uma magia de Cura de rank Imperador --- o topo da escada, engarrafado. \\
\end{tabelalivro}

\begin{tabelalivro}{L{4.8cm} C{1.6cm} C{1.4cm} L{5.6cm}}{Venenos --- Rank de Guilda é o mínimo pra o item aparecer à venda}
\cabtabela{\textbf{Item} & \textbf{Preço} & \textbf{Rank} & \textbf{Notas}}
Veneno Fraco (Profundidade 1) & 5 PO & E & Ex: baba de sapo-lodo. Aplicação em Cap. 4, §7, "Aplicando um Veneno em Combate ou em Segredo". \\
Veneno Comum (Profundidade 2) & 20 PO & C & Ex: peçonha de serpente-do-pântano. Venda exige licença registrada na sede. \\
Veneno Potente (Profundidade 3) & 80 PO & A & Ex: fel de wyvern. Venda sob vigilância da sede --- vender sem licença é crime em Millis e no Reino Asura (perde Reputação com a facção local). \\
Kit de Coleta de Veneno & 15 PO & F & Frascos, pinça e luvas grossas --- permite extrair a glândula ou a baba de uma criatura recém-abatida antes que o veneno perca a potência. \\
Kit de Aplicação Segura & 40 PO & D & Aplicador de cabo longo e antídoto de emergência incluso --- reduz (não elimina) o risco de se envenenar aplicando a própria dose. \\
Frasco de Armazenamento Estável & 100 PO & B & Vidro escuro selado a vácuo --- uma dose guardada aqui não perde Profundidade com o tempo, ao contrário de um frasco comum. \\
Antídoto Universal & 500 PO & S & Remove TODA a Profundidade acumulada de uma única aflição (Cap. 4, §7), não importa o quão fundo --- inclusive Profundidade 4+, que nenhuma dose à venda alcança. Raríssimo por isso mesmo. \\
\end{tabelalivro}

\begin{tabelalivro}{L{4.8cm} C{1.6cm} C{1.4cm} L{5.6cm}}{Ferramentas Mágicas --- Rank de Guilda é o mínimo pra o item aparecer à venda}
\cabtabela{\textbf{Item} & \textbf{Preço} & \textbf{Rank} & \textbf{Notas}}
Pedra de Aquecimento & 25 PO & F & Esquenta ao segurar com um fio de mana. Comum em acampamento de inverno. \\
Caixa Resfriadora & 60 PO & E & Mantém comida e bebida fria por dias sem gelo --- item de conforto, não de combate. \\
Pedra de Detecção de Mana & 150 PO & D & Brilha perto de fontes de mana ativa num raio curto --- acha magos escondidos ou armadilhas encantadas, mas não distingue amigo de inimigo. \\
Pedra de Comunicação (par) & 220 PO & D & Duas pedras gravadas juntas: o que se escreve na principal aparece na secundária, não importa a distância. Só texto --- não substitui magia. \\
Lançador de Pergaminhos & 400 PO & B & Suporte dorsal com dez encaixes pra pergaminhos de magia pré-preparados --- um toque de mana ativa o encaixe certo sem precisar sacar nada com as mãos. Cada pergaminho é preparado à parte (gancho de campanha, não item de prateleira). \\
Capa Termorreguladora & 45 PO & E & Mantém sua temperatura corporal estável em qualquer clima comum --- não protege contra frio ou calor de origem mágica. \\
Luvas Amortecedoras & 70 PO & D & Absorvem metade do impacto que a própria mão receberia ao golpear algo sólido --- poupa os nós dos dedos de quem luta desarmado. \\
Óculos de Detecção & 160 PO & D & Realçam o contorno de quem se esconde ou se disfarça, num raio curto --- não atravessa Invisibilidade mágica de verdade. \\
Botas Velozes & 180 PO & C & Encantadas pra dobrar sua velocidade de corrida por curtos períodos --- o Mestre define duração e limite de uso. Não empilha com outros efeitos de deslocamento. \\
Cajado de Aprendiz & 12 PO & F & Ajuda a canalizar mana com mais estilo que as próprias mãos --- efeito puramente estético, sem bônus mecânico. \\
Brinco de Ilusão Breve & 350 PO & B & Cravado na pele com 1 Ação, cria uma distração visual ou muda levemente sua aparência por até 1 cena --- duração e limites exatos ficam a critério do Mestre. \\
Anel Perturbador de Magia & 900 PO & S & Apontado pra um alvo, faz a magia dele simplesmente falhar ao tentar operar perto de você por alguns instantes. Raríssimo --- como o Anel de Teleporte (Cap. 5 §4), cabe ao Mestre decidir se um artefato desse nível existe na sua campanha. \\
Anel de Barreira Improvisada & 900 PO & S & Apontado a tempo, ergue uma barreira invisível que intercepta um único golpe físico vindo contra você. Mesmo aviso do item acima: raro, e o Mestre decide se existe na mesa. \\
Bússola Encantada & 90 PO & E & O ponteiro sempre aponta pra um lugar gravado nela (a sede da Guilda mais próxima, por padrão) em vez do norte --- reduz Desvantagem por estar perdido, mas não substitui um mapa. \\
Amuleto de Respiração Aquática & 140 PO & D & Enquanto usado, respira embaixo d'água normalmente por um tempo curto --- o Mestre define a duração. Não ajuda contra pressão de profundidade nem frio. \\
Amuleto de Resistência ao Fogo & 500 PO & A & Enquanto usado, reduz o desconforto de calor extremo (deserto, vulcão, fornalha) a um incômodo administrável. Não é a Resistência a dano de fogo em combate (Cap. 4) --- isso continua exigindo magia ou talento de verdade. \\
Sino de Alarme Mágico & 55 PO & E & Fincado no chão, toca sozinho (só audível a quem o ativou) se alguém cruzar um perímetro curto ao redor --- clássico de acampamento em território hostil. \\
\end{tabelalivro}

\begin{tabelalivro}{L{4.8cm} C{1.6cm} C{1.4cm} L{5.6cm}}{Encantamentos --- Rank de Guilda é o mínimo pra o item aparecer à venda}
\cabtabela{\textbf{Item} & \textbf{Preço} & \textbf{Rank} & \textbf{Notas}}
Encantamento Avançado & 150 PO & D & Serviço encomendado: +1 no Dado de Arma OU +1 na CA de um item que você já possui. Encantador precisa de Rank Avançado na árvore compatível. Depois de comprar, edite o item alvo pra refletir o bônus. \\
Encantamento Santo & 300 PO & B & Serviço encomendado: dano elemental extra (+1d6, tipo à escolha) num item que você já possui. Encantador precisa de Rank Santo na árvore compatível. Depois de comprar, edite o item alvo pra refletir o bônus. \\
Encantamento Rei & 600 PO & A & Serviço encomendado: ignora Resistência a um tipo de dano num item que você já possui. Encantador precisa de Rank Rei na árvore compatível. Depois de comprar, edite o item alvo pra refletir o bônus. \\
Encantamento Imperador & 1500 PO & S & Serviço encomendado: +1 no Bônus de Rank pra fins de Dado de Arma, num item que você já possui. Encantador precisa de Rank Imperador na árvore compatível. Depois de comprar, edite o item alvo pra refletir o bônus. \\
\end{tabelalivro}



% ---------------------------------------------------------------------------
\section{Reputação com Facções}

Nem toda consequência cabe em PA ou em Rank de Aventureiro --- às vezes o que muda é quem
abre a porta pra você. Reputação é uma escala narrativa de cinco degraus, uma por facção,
movida pelo Mestre conforme os atos públicos do grupo.

\begin{nota}{Guilda ou Reputação --- qual eu uso?}
    As duas medem \enquote{quão bem o mundo te trata}, mas em escalas diferentes: a
    \textbf{Guilda} é uma fama profissional única, a mesma em qualquer cidade que tenha uma
    sede. \textbf{Reputação} é por facção específica --- dá pra ser Aliado do Reino Asura e
    Inimigo da Igreja de Millis ao mesmo tempo, e nenhuma das duas mexe na outra.
\end{nota}

\begin{tabelalivro}{L{3cm} L{10.8cm}}{Os cinco degraus}
\cabtabela{\textbf{Nível} & \textbf{O que significa, em qualquer facção}}
Inimigo ($-2$) & A facção age ativamente contra o grupo, sempre que puder fazer isso sem custo alto pra ela. \\
Desconfiado ($-1$) & Portas se fecham por precaução. Nenhum ataque direto, mas nenhuma ajuda também. \\
Neutro (0) & Ponto de partida padrão --- a facção nem sabe quem vocês são. \\
Respeitado ($+1$) & Contratos, favores e informação ficam mais fáceis dentro do território da facção. \\
Aliado ($+2$) & A facção arrisca recursos reais pelo grupo --- tropas, magos, dinheiro. \\
\end{tabelalivro}

\begin{tabelalivro}{L{2.4cm} L{3.8cm} L{3.8cm} L{3.8cm}}{As Três Facções deste Livro}
\cabtabela{\textbf{Nível} & \textbf{Reino Asura} & \textbf{Igreja de Millis} & \textbf{Deuses Demônios}}
Inimigo & Mandado de captura ativo --- a guarda ataca de vista. & Excomungado. Templos recusam cura, abrigo e água. & Marcado como inimigo pela Imperatriz Kishirika. \\
Desconfiado & Vigiados: espiões da coroa relatam cada movimento em Ars. & Sacerdotes recusam bênção e informação, mas não interferem. & Tolerados, desde que fiquem fora do território de um clã específico. \\
Neutro & Só mais um grupo no registro da capital. & Nenhum templo conhece o grupo pelo nome. & O grupo é estrangeiro --- cuidado padrão, nada pessoal. \\
Respeitado & Acesso à corte menor; contratos diretos, sem passar pela Guilda. & Curas gratuitas em templos menores; acesso à biblioteca de um mosteiro. & Um clã garante passagem segura pelo território. \\
Aliado & Audiência com a coroa por pedido; tropas reais em campanhas regionais. & O Grande Templo de Millis abre arquivos restritos. & A Imperatriz Kishirika reconhece o grupo. \\
\end{tabelalivro}

\begin{nota}{Como o Mestre move o marcador}
    Não existe fórmula. Reputação sobe ou desce por atos públicos, não por PA gasto ou sessões
    jogadas. Mude só um degrau por vez, e só quando o ato for grande o bastante pra virar
    boato ou registro oficial. E facções não são unânimes: nada impede um personagem de ser
    Aliado de uma ala e Inimigo de outra dentro da mesma facção nominal.
\end{nota}

% ---------------------------------------------------------------------------
\section{Crafting e Alquimia}

Quatro perguntas resolvem qualquer fabricação: quem pode fazer, quanto tempo leva, quanto
custa em materiais e o que acontece se o teste falhar.

\begin{itemize}
    \item \textbf{Quem:} qualquer personagem com a Perícia de Ofícios (especializada em
          Alquimia, pra Poções e Venenos) ligada ao item. Encantamento é diferente --- exige
          um encantador vivo no Rank listado, não a Perícia de Ofícios.
    \item \textbf{Até onde:} uma poção que reproduz magia de rank X exige que o alquimista
          tenha um patamar igual ou superior a X em Cura (ou em Desintoxicação, pros
          antídotos). Um veneno de Profundidade X exige um patamar igual ou superior a X em
          Desintoxicação, ou material colhido de uma criatura daquele porte. Sem isso, a
          receita simplesmente não é legível --- não é uma CD mais alta, é um teste que você
          não pode tentar.
    \item \textbf{Tempo:} 1 bloco de Downtime por item, salvo quando a tabela disser outro
          valor.
    \item \textbf{Custo em materiais:} metade do valor listado na coluna de Custo. O valor
          cheio é o preço de venda, não o de fabricação.
    \item \textbf{Teste:} role Ofícios contra a CD da tabela ao fim do bloco. Sucesso: o item
          fica pronto. Falha: os materiais se perdem, mas o bloco já foi gasto. Falha crítica
          (1 no dado): metade dos materiais é recuperável.
\end{itemize}

\begin{aviso}{Por que existe um portão de Rank, e não só uma CD alta}
    Sem ele, a CD era o único obstáculo --- e a Perícia de Ofícios dá Vantagem (2d20). Um
    alquimista de 1º patamar com Intelecto 4 batia a CD 20 da Poção Imperial de Cura em pouco
    mais da metade das tentativas, e um grupo com uma semana livre por arco engarrafava cura
    de rank Imperador antes de conhecer um mago de rank Avançado. O portão da Guilda cobria só
    a \emph{compra}; a fabricação passava por baixo dele.

    \medskip
    É a mesma lógica que o Encantamento já usava --- lá o portão sempre foi o Rank do
    encantador, nunca uma rolagem. Agora as duas metades desta seção cobram a mesma coisa:
    você fabrica o que você entende, e paga em PO o que você não entende.
\end{aviso}

\subsection{Poções}

\begin{tabelalivro}{L{4.4cm} C{1.4cm} C{2.6cm} L{5.4cm}}{}
\cabtabela{\textbf{Poção} & \textbf{CD} & \textbf{Venda / Fabricação} & \textbf{Efeito}}
Poção Menor de Cura & 11 & 15 / 8 PO & Reproduz uma magia de Cura de rank Principiante ou Intermediário. \\
Poção de Antídoto & 13 & 25 / 13 PO & Remove 1 ponto de Profundidade de uma única aflição. \\
Poção Maior de Cura & 15 & 60 / 30 PO & Reproduz uma magia de Cura de rank Avançado ou Santo. \\
Elixir de Foco & 13 & 40 / 20 PO & Vantagem no próximo teste de resistência de Espírito. \\
Poção de Vigor Passageiro & 14 & 45 / 23 PO & Vantagem no próximo teste de resistência de Vigor. \\
Poção Régia de Cura & 17 & 120 / 60 PO & Reproduz uma magia de Cura de rank Rei. \\
Elixir de Regeneração & 18 & 200 / 100 PO & Remove 2 níveis de Exaustão. Não remove Exaustão cuja causa ainda esteja ativa: quem não comeu continua com fome. \\
Poção Imperial de Cura & 20 & 400 / 200 PO & Reproduz uma magia de Cura de rank Imperador --- o topo da escada, engarrafado. \\
\end{tabelalivro}

\subsection{Venenos}

Fabricar veneno é produzir uma dose de uma aflição já catalogada no Cap. 4, §7 --- a CD sobe
junto com a Profundidade, porque manusear algo mais perigoso sem se envenenar no processo é
mais difícil.

\begin{tabelalivro}{C{2.4cm} L{5.4cm} C{1.4cm} L{4.4cm}}{}
\cabtabela{\textbf{Profundidade} & \textbf{Exemplo} & \textbf{CD} & \textbf{Venda / Fabricação}}
1 & Baba de Sapo-Lodo & 10 & 5 / 3 PO \\
2 & Peçonha de Serpente-do-Pântano & 12 & 20 / 10 PO \\
3 & Fel de Wyvern & 14 & 80 / 40 PO \\
4$+$ & Praga do Continente Demônio & 16$+$ & Não está à venda --- só se rouba, caça ou herda. \\
\end{tabelalivro}

\begin{aviso}{A lei e o veneno}
    Vender veneno de Profundidade 3 ou superior sem licença é crime em Millis e no Reino
    Asura --- perde Reputação com a facção local automaticamente.
\end{aviso}

\subsection{Encantamento de Arma e Armadura}

Encantar não é uma Perícia de Ofícios --- é um serviço prestado por um mago que já alcançou o
Rank exigido numa árvore compatível com o efeito (dano elemental pede a Magia daquele
elemento; resistência e CA pedem Barreira; qualquer efeito genérico aceita Invocação).

\begin{tabelalivro}{L{6.4cm} C{2.4cm} C{1.8cm} C{2.2cm}}{}
\cabtabela{\textbf{Efeito} & \textbf{Rank exigido} & \textbf{Tempo} & \textbf{Custo}}
$+1$ no Dado de Arma ou $+1$ na CA & Avançado & 1 bloco & 150 PO \\
Dano elemental extra ($+$1d6, tipo à escolha) & Santo & 2 blocos & 300 PO \\
Ignora Resistência a um tipo de dano & Rei & 4 blocos & 600 PO \\
$+1$ degrau na Escada de Dados de Arma & Imperador & 8 blocos & 1500 PO \\
\end{tabelalivro}

\begin{itemize}
    \item Um item só carrega um encantamento por vez. Encantar de novo substitui o anterior.
          O encantamento de Imperador é a única exceção parcial: ele sobe um degrau da
          Escada, e um item que já tinha o $+1$ do Avançado sobe dois no total.
    \item O custo em PO já é o total (materiais $+$ trabalho) --- não se aplica a divisão por
          metade do Downtime comum, porque não é o personagem fazendo o trabalho manual.
    \item O item-base precisa estar em posse do encantador durante todo o tempo listado ---
          ele não trabalha à distância.
    \item Não existe teste de falha aqui: se o encantador tem o Rank exigido, tempo e PO
          cobrem o serviço inteiro.
\end{itemize}

\subsection{Itens Mágicos Únicos}

Nem todo item mágico cabe numa tabela de preço. Alguns são artefatos: peças únicas cuja
fabricação é evento de campanha, não compra de ficha --- exatamente como o Rank Deus.

\begin{exemplo}{O Anel de Teleporte}
    \textbf{O que faz:} teleporta o portador --- e quem ele tocar --- pra um de até três
    destinos gravados, sem custo de PM, sem teste. Depois de usado, precisa de 1 semana pra
    recarregar.

    \medskip
    \textbf{Por que não está na tabela acima:} exige um encantador de Invocação de rank Deus,
    materiais que só existem em circunstâncias específicas da campanha, e é irrepetível ---
    gravar um novo destino exige voltar ao mesmo encantador. Trate a fabricação como o final
    de um arco inteiro.
\end{exemplo}
